\setupcolors[state=start]
\setupcolor[hex]
\definecolor[darkblue][h=003366]
\definecolor[darkgray][h=e6e6e6]
\definecolor[redline][h=d9534f]
\definecolor[blueline][h=5bc0de]
\definecolor[greenline][h=5cb85c]

\setuphead[section][style=\bfa, color=darkblue, before=\blank, after=\blank]
\setuphead[subsection][style=\bf, color=black, before=\blank, after=\blank]
\setuppapersize[A4][A4]
\setupbodyfont[11pt]
\setuplayout[width=middle, height=middle, topspace=2cm, bottomspace=2cm]
\mainlanguage[es]

\starttext
\title{\bf Comparación de OPP\_v5\_1\_mono\crlf y OPP\_v5\_2\_split en pruebas funcionales y de stress}
\blank[big]

\startsection[title={Objetivo y alcance}]
Este documento compara las ejecuciones de los tests funcionales y de stress de OPP\_v5\_1\_mono
(formulación monolítica) y OPP\_v5\_2\_split (formulación en dos etapas) sobre los mismos
escenarios. Los dos modelos definen el mismo problema de optimización (ruteo, asignación y
secuenciación, con velocidades heterogéneas de pickers y cadena de frío) y tienen el mismo valor
óptimo.

Se usan las corridas del 24 de septiembre de 2026, que ya incluyen los tests 35 a 38 (cadena de
frío, velocidad, demora a propósito y el límite de un arco por vértice):
\startitemize
\item \type{results/opp_v5_2_split/functional/} — %%NUM:split_tests%% tests, corrida sola.
\item \type{results/opp_v5_1_mono/functional/} — %%NUM:mono_tests%% tests, corrida mientras se
ejecutaba aparte el test\_case\_30 de los dos modelos, así que sus tiempos pueden estar algo
inflados.
\stopitemize

El test\_case\_30 se corrió por separado y sus filas se agregaron a cada CSV. Su variante
factible no termina en los 25 minutos de tope en ninguno de los dos modelos y queda fuera de la
estadística; la infactible figura con los demás infactibles.

La suite de stress (sección \in[sec:stress]) usa \type{results/opp_v5_2_split/stress/} y
\type{results/opp_v5_1_mono/stress/}: los mismos %%NUM:stress_total%% escenarios en los dos
modelos, con un tope de 40 minutos por test. Los estudios que explican esos resultados están en
los documentos de exploración de cada modelo.
\stopsection

\startsection[title={Criterio de equivalencia}]
Los dos modelos resuelven la misma definición: misma región factible, misma función objetivo,
mismo valor óptimo. El argumento es que la ruta más corta domina a cualquier otra (el objetivo no
tiene un término de tiempo de ruta, y una ruta más corta solo afloja las restricciones de
duración y de cadena de frío), y los lotes no interactúan en el ruteo.

Lo que difiere es el {\em conjunto} de soluciones óptimas: el monolítico puede devolver una ruta
no mínima con el mismo valor objetivo. En el split, la Etapa 1 siempre produce la ruta mínima,
porque su objetivo es minimizar el tiempo de ruta.

Los tests de los dos modelos se corresponden uno a uno, salvo
\hyphenatedurl{test_preflight_detects_when_no_picker_is_eligible} (test\_case\_24): resuelve solo la
Etapa 1, que el monolítico no tiene. Los escenarios heredados de OPP\_v4 neutralizan las
capacidades de v5 con $V_p = 1$ y un límite de frío que no restringe.
\stopsection

\startsection[title={Cobertura y resultado de la ejecución}]
%%TABLE:coverage%%

Los infactibles incluyen los tests 34 y 35 de cadena de frío. En el CSV del split figuran con un
objetivo porque su chequeo previo corta antes de la Etapa 2 y queda registrado el objetivo de la
Etapa 1; los dos modelos los declaran infactibles.
\stopsection

\startsection[title={Equivalencia del valor objetivo}]
De los %%NUM:feasible%% pares factibles, {\bf %%NUM:same_obj%% tienen el mismo valor objetivo}. La
diferencia máxima es $%%NUM:max_diff%%$ (casos %%NUM:diff_cases%%), dentro de la tolerancia de
SCIP.

Esto confirma empíricamente que las dos formulaciones tienen el mismo valor óptimo, también en
los tests nuevos de cadena de frío, velocidad y demora a propósito.
\stopsection

\startsection[title={Tiempos de los pares factibles}]
Sobre los %%NUM:feasible%% pares factibles. \type{total_time} es la ejecución completa del test,
incluida la compilación de ZIMPL; \type{model_time} es solo el solver (en el split, la Etapa 2).

%%TABLE:times%%

En total, el monolítico tardó %%NUM:sum_ratio%% veces más que el split, pero en la mediana la
diferencia se invierte: %%NUM:median_mono%%~s contra %%NUM:median_split%%~s. En %%NUM:mono_faster%%
de los %%NUM:feasible%% pares el monolítico fue {\bf más rápido} (entre %%NUM:mono_faster_min%%× y
%%NUM:mono_faster_max%%×): son instancias chicas, donde compilar y resolver un solo modelo cuesta
menos que compilar dos. La suma la dominan las instancias de grilla, donde el monolítico llega a
tardar %%NUM:top_ratio%% veces más (caso %%NUM:top_case%%).

En \type{model_time} la diferencia es mayor (%%NUM:model_sum_ratio%% veces en la suma,
%%NUM:model_median_ratio%% en la mediana), porque el monolítico resuelve un MIP con más variables y
restricciones que cualquiera de las dos etapas.
\stopsection

\startsection[title={Visualización comparativa}]

\startsubsection[title={Tiempo total: monolítico contra dos etapas}]
Cada punto es un par factible. Los puntos por encima de la diagonal son instancias donde el
monolítico fue más lento.

\placefigure[force][fig:scatter]{Tiempo total de cada par factible: monolítico contra split (escala logarítmica en los dos ejes).}
{%%FIG:scatter%%}
\stopsubsection

\startsubsection[title={Cociente de tiempos por par}]
Barras ordenadas por el cociente mono / split. Las que quedan por debajo de 1 son instancias
donde el monolítico fue más rápido. Se identifican los casos en que el monolítico tarda más de
10 veces.

\placefigure[force][fig:ratio]{Tiempo total del monolítico dividido el del split, por par factible (escala logarítmica).}
{%%FIG:ratio%%}
\stopsubsection

\startsubsection[title={Tiempo del solver según el tamaño}]
Tiempo del solver contra el tamaño de la instancia, medido como lotes × ubicaciones, para los
%%NUM:sized%% pares cuyo CSV registra ese tamaño. La separación entre las dos nubes crece con el
tamaño: la descomposición escala mejor.

\placefigure[force][fig:size]{Tiempo del solver contra el tamaño de la instancia (escala logarítmica en los dos ejes).}
{%%FIG:size%%}
\stopsubsection
\stopsection

\startsection[title={Detalle de los pares factibles}]
%%TABLE:detail%%
\stopsection

\startsection[title={Tests infactibles}]
Los %%NUM:infeasible%% pares que los dos modelos declaran infactibles.

%%TABLE:infeasible%%
\stopsection

\startsection[title={Tests fuera de la estadística}]
%%TABLE:excluded%%
\stopsection

\startsection[title={Análisis de la falla del test\_case\_06}]
El escenario del test\_case\_06 (\type{test_two_possible_routes}) tiene dos rutas que producen el
mismo valor objetivo (30,0): L1→L2→L4→L5, con un tiempo de ruta de 7,5, y L1→L4→L5, con 4,0. Las
dos son factibles y alcanzan el mismo objetivo porque la fecha de entrega de 15,0 da margen para
cualquiera.

El split siempre devuelve la ruta mínima (4,0), porque la Etapa 1 minimiza el tiempo de ruta. El
monolítico puede elegir la más larga (7,5), porque su objetivo no incluye el tiempo de ruta y las
dos rutas dan el mismo valor.

El test falla porque su aserción verifica el grafo de arcos de la ruta mínima. El monolítico no
está en error: devuelve una solución óptima válida. Es una muestra de que los conjuntos de
soluciones óptimas difieren entre las formulaciones. La falla está registrada como
\type{known_failures} en \type{scripts/baseline_counts.json}.
\stopsection

\startsection[title={Suite de stress},reference=sec:stress]
Los tests de stress llevan los dos modelos a instancias de tamaño real: el perfil realista (26
órdenes por lote, dos bandas de velocidad, pesos 4/5) en grillas de 7×7 a 16×16 y cuadrillas de
hasta 24 lotes, más tres escenarios heredados de OPP\_v4. El split corre la Etapa 1 con seis
procesos. Su tiempo es la duración completa del test y es un techo: la máquina entró en
suspensión durante parte de esa corrida, y el tiempo suspendido cuenta en la duración pero no en
el tope.

%%TABLE:stress%%

El monolítico termina %%NUM:stress_mono_ok%% de los %%NUM:stress_total%% tests; el split,
%%NUM:stress_split_ok%%, ninguno en más de %%NUM:stress_split_max%% minutos. En los tests
funcionales las dos formulaciones son intercambiables salvo en las grillas; a este tamaño, solo
el split da una respuesta. Los tests que tampoco termina el split son los del perfil realista
con 8 lotes o más (12 a 14 y 16 a 19), la grilla de 16×16 (08) y dos de los heredados (01 y
03).
\stopsection

\startsection[title={Conclusiones}]
\startitemize[n]
\item {\bf Equivalencia del valor objetivo:} confirmada en los %%NUM:feasible%% pares factibles,
incluidos los tests nuevos. La diferencia máxima, $%%NUM:max_diff%%$, es de redondeo del solver.
\item {\bf Instancias chicas:} el monolítico es comparable o algo más rápido, porque evita
compilar dos modelos ZIMPL; el split paga un costo fijo por la segunda etapa.
\item {\bf Instancias de grilla:} el monolítico es entre 10 y %%NUM:top_ratio%% veces más lento. El
MIP combinado tiene bastantes más variables y restricciones que cualquiera de las dos etapas.
\item {\bf Diferencia de conjuntos de soluciones:} el test\_case\_06 muestra que el monolítico
puede devolver rutas no mínimas con el mismo objetivo.
\item {\bf Escalabilidad:} la ventaja del split crece con el tamaño de la instancia. En la suite
de stress el monolítico termina %%NUM:stress_mono_ok%% de %%NUM:stress_total%% tests en 40
minutos y el split %%NUM:stress_split_ok%%.
\stopitemize
\stopsection

\stoptext
